\documentclass[11pt]{article}
\usepackage[margin=1in]{geometry}
\usepackage{amsmath,amssymb}
\usepackage{enumitem}
\usepackage{booktabs}
\usepackage{hyperref}

\title{\textbf{Course Structure \& Detailed Breakdown}\\ \large Math 617: Number Theory I (Johns Hopkins University)}
\author{Based on Syllabus Data}
\date{Fall Semester}

\begin{document}

\maketitle

\section*{I. Course Overview \& Architecture}

\begin{itemize}[leftmargin=*]
    \item \textbf{Level:} Graduate level (Requires mastery of qualifying exam-level Algebra and Analysis).
    \item \textbf{Assessment Weight:} $100\%$ Homework (Weekly problem sets; lowest score dropped).
    \item \textbf{Primary References:} 
        \begin{itemize}
            \item Jürgen Neukirch, \emph{Algebraic Number Theory}
            \item André Weil, \emph{Basic Number Theory}
            \item Serge Lang, \emph{Algebraic Number Theory}
            \item J.W.S. Cassels \& A. Fröhlich (Eds.), \emph{Algebraic Number Theory}
        \end{itemize}
\end{itemize}

---

\section*{II. Core Prerequisite Knowledge}

To succeed in this course, students must have firm ground in the following topics:

\begin{description}[leftmargin=*]
    \item[Graduate Algebra:] Ring theory (PIDs, UFDs, maximal/prime ideals), Module theory (structure theorem for finitely generated modules over PIDs), Galois theory (field extensions, Galois groups, trace, and norm maps), and group representation basics.
    \item[Graduate Analysis:] Measure theory and integration (Haar measure on locally compact abelian groups), basic topology (topological groups, compactness, completion of metric spaces), and complex analysis (analytic continuation, gamma function basics).
\end{description}

---

\section*{III. Modular Course Breakdown}

\subsection*{Module 1: Dedekind Domains and Ramification Theory}
\textbf{Focus:} Structural algebraic number theory; moving from ideal factorization in rings of integers to geometric and local behavior via ramification.

\begin{itemize}[leftmargin=*]
    \item \textbf{Key Concepts:}
        \begin{itemize}
            \item Noetherian, integrally closed domains of Krull dimension 1.
            \item Unique factorization of fractional ideals.
            \item Prime ideal decomposition in extension fields: $e$ (ramification index), $f$ (inertia degree), and $g$ (number of prime factors). The fundamental identity $\sum_{i=1}^g e_i f_i = n$.
            \item Galois action on prime ideals, decomposition groups ($D_\mathfrak{P}$), and inertia groups ($I_\mathfrak{P}$).
            \item Discriminant and Different: Different ideal $\mathfrak{D}_{L/K}$, discriminant $\mathfrak{d}_{L/K}$, and Dedekind's criterion for ramification.
        \end{itemize}
    \item \textbf{Learning Objectives:}
        \begin{itemize}
            \item Prove the unique factorization of ideals in Dedekind domains.
            \item Compute ramification indices and residue degree extensions in explicit field extensions (e.g., cyclotomic and quadratic fields).
            \item Analyze unramified, tamely ramified, and wildly ramified extensions via lower and upper filtration series of ramification groups.
        \end{itemize}
    \item \textbf{Identified Difficult Topics:} Distinguishing between wild and tame ramification; precise calculation of the relative different and its connection to algebraic ramification.
\end{itemize}

\subsection*{Module 2: Global and Local Fields}
\textbf{Focus:} Transitioning from global field arithmetic (number fields, function fields over finite fields) to local fields via valuations and completions ($p$-adic numbers).

\begin{itemize}[leftmargin=*]
    \item \textbf{Key Concepts:}
        \begin{itemize}
            \item Absolute values, valuations (archimedean vs. non-archimedean), and Ostrowski's Theorem.
            \item Completions of global fields: $p$-adic numbers $\mathbb{Q}_p$ and local fields $K_v$.
            \item Structure of local fields: Ring of integers $\mathcal{O}_v$, maximal ideal $\mathfrak{m}_v$, residue field $k_v$.
            \item Hensel's Lemma (multiple formulations) and roots of unity in $\mathbb{Q}_p$.
            \item Unramified and totally ramified extensions of local fields.
            \item Local norm-index computations.
        \end{itemize}
    \item \textbf{Learning Objectives:}
        \begin{itemize}
            \item Classify all non-trivial topologies/absolute values on $\mathbb{Q}$ and global function fields.
            \item Apply Hensel's Lemma to solve polynomial equations in local fields.
            \item Formulate the local-global principle (Hasse Principle) for quadratic forms.
        \end{itemize}
    \item \textbf{Identified Difficult Topics:} Local-to-global spectral arguments; topological groups over non-archimedean fields; local class field theory inputs.
\end{itemize}

\subsection*{Module 3: Adèles and Idèles}
\textbf{Focus:} Unifying local extensions into global topological objects to examine global number theory through harmonic analysis.

\begin{itemize}[leftmargin=*]
    \item \textbf{Key Concepts:}
        \begin{itemize}
            \item Restricted direct products of topological spaces.
            \item The Ring of Adèles $\mathbb{A}_K$: Topology, compactness of $\mathbb{A}_K / K$, and discrete embeddings.
            \item The Group of Idèles $\mathbb{I}_K$: Idele class group $C_K = \mathbb{I}_K / K^\times$, compactness of $C_K^1$.
            \item Connection between idele classes, ideal class groups, and Dirichlet's Unit Theorem.
        \end{itemize}
    \item \textbf{Learning Objectives:}
        \begin{itemize}
            \item Construct the restricted product topology and prove structural properties of $\mathbb{A}_K$.
            \item Show that $K$ is discrete and co-compact in $\mathbb{A}_K$.
            \item Translate classical ideal theory statements into the language of adèles/idèles.
        \end{itemize}
    \item \textbf{Identified Difficult Topics:} Working fluently with the restricted product topology; establishing co-compactness of $K \subset \mathbb{A}_K$.
\end{itemize}

\subsection*{Module 4: Tate's Thesis (Harmonic Analysis on Local/Global Fields)}
\textbf{Focus:} Analytic continuation and functional equations of global $L$-functions via Fourier analysis on adèles (if time permits).

\begin{itemize}[leftmargin=*]
    \item \textbf{Key Concepts:}
        \begin{itemize}
            \item Haar measures on local fields and adèle rings.
            \item Local Fourier transforms, self-duality, and local gamma factors.
            \item Schwartz-Bruhat functions on adèles.
            \item Tate's global zeta integral $Z(f, \chi)$ and the Poisson summation formula on $\mathbb{A}_K$.
            \item Functional equation and analytic continuation of Riemann Zeta and Dirichlet $L$-functions.
        \end{itemize}
    \item \textbf{Learning Objectives:}
        \begin{itemize}
            \item Calculate local zeta integrals for unramified and ramified characters.
            \item Use Poisson summation on the adèle space to derive global functional equations cleanly.
        \end{itemize}
    \item \textbf{Identified Difficult Topics:} Harmonizing local and global Fourier transforms; subtle topology of characters on the idèle class group.
\end{itemize}

---

\section*{IV. Summary Map of Pedagogical Dependencies}

\begin{table}[h!]
\centering
\small
\begin{tabular}{@{}lll@{}}
\toprule
\textbf{Module} & \textbf{Core Prerequisites} & \textbf{Conceptual Peak / Goal} \\ \midrule
Module 1: Dedekind Domains & Abstract Algebra (Galois, Ideal Theory) & Ramification Groups \& Different Ideal \\
Module 2: Local Fields & Point-Set Topology, Analysis & Hensel's Lemma \& Valuations \\
Module 3: Adèles \& Idèles & Modules 1 \& 2, Topological Groups & Compactness of $C_K^1$ \\
Module 4: Tate's Thesis & Module 3, Harmonic Analysis & Functional Equation of $L$-functions \\ \bottomrule
\end{tabular}
\caption{Progression of mathematical tools across the semester.}
\end{table}

\end{document}